% math0 manuscript. Build with: latexmk -pdf paper.tex
\documentclass[11pt]{amsart}
\usepackage[margin=1.1in]{geometry}
\usepackage{amsmath,amssymb,amsthm,mathtools}
\usepackage{url}
\usepackage{hyperref}
\usepackage{xcolor}
\usepackage{listings}
\usepackage{newunicodechar}

\lstdefinelanguage{lean}{
  keywords={theorem,lemma,def,abbrev,example,fun,by,have,show,from,let,in,if,then,else,match,with,
    namespace,end,open,import,section,variable,instance,structure,where,calc,obtain,intro,exact,
    refine,rcases,rintro,simp,omega,decide,norm_num,linarith,nlinarith,positivity,induction,cases,
    private,rw},
  sensitive=true, comment=[l]{--}, morecomment=[s]{/-}{-/}, string=[b]"}
\lstset{language=lean, basicstyle=\ttfamily\footnotesize, columns=fullflexible, keepspaces=true, breaklines=true,
  keywordstyle=\color{blue!60!black}\bfseries, commentstyle=\color{gray}, frame=single,
  extendedchars=true, inputencoding=utf8,
  literate={ℕ}{{$\mathbb{N}$}}1 {∀}{{$\forall$}}1 {∃}{{$\exists$}}1 {¬}{{$\neg$}}1 {∧}{{$\wedge$}}1
    {→}{{$\to$}}1 {≤}{{$\le$}}1 {∈}{{$\in$}}1 {·}{{$\cdot$}}1}
% Unicode commonly found in Lean sources.
\newunicodechar{ℕ}{\ensuremath{\mathbb{N}}}
\newunicodechar{ℤ}{\ensuremath{\mathbb{Z}}}
\newunicodechar{ℚ}{\ensuremath{\mathbb{Q}}}
\newunicodechar{ℝ}{\ensuremath{\mathbb{R}}}
\newunicodechar{ℂ}{\ensuremath{\mathbb{C}}}
\newunicodechar{∀}{\ensuremath{\forall}}
\newunicodechar{∃}{\ensuremath{\exists}}
\newunicodechar{→}{\ensuremath{\to}}
\newunicodechar{←}{\ensuremath{\leftarrow}}
\newunicodechar{↔}{\ensuremath{\leftrightarrow}}
\newunicodechar{↦}{\ensuremath{\mapsto}}
\newunicodechar{¬}{\ensuremath{\neg}}
\newunicodechar{∧}{\ensuremath{\wedge}}
\newunicodechar{∨}{\ensuremath{\vee}}
\newunicodechar{≤}{\ensuremath{\le}}
\newunicodechar{≥}{\ensuremath{\ge}}
\newunicodechar{≠}{\ensuremath{\ne}}
\newunicodechar{∈}{\ensuremath{\in}}
\newunicodechar{∉}{\ensuremath{\notin}}
\newunicodechar{⊆}{\ensuremath{\subseteq}}
\newunicodechar{∩}{\ensuremath{\cap}}
\newunicodechar{∪}{\ensuremath{\cup}}
\newunicodechar{∑}{\ensuremath{\sum}}
\newunicodechar{∏}{\ensuremath{\prod}}
\newunicodechar{⟨}{\ensuremath{\langle}}
\newunicodechar{⟩}{\ensuremath{\rangle}}
\newunicodechar{·}{\ensuremath{\cdot}}
\newunicodechar{∘}{\ensuremath{\circ}}
\newunicodechar{⁻}{\ensuremath{^{-}}}
\newunicodechar{¹}{\ensuremath{^{1}}}
\newunicodechar{²}{\ensuremath{^{2}}}
\newunicodechar{₀}{\ensuremath{_{0}}}
\newunicodechar{₁}{\ensuremath{_{1}}}
\newunicodechar{₂}{\ensuremath{_{2}}}
\newunicodechar{α}{\ensuremath{\alpha}}
\newunicodechar{β}{\ensuremath{\beta}}
\newunicodechar{ε}{\ensuremath{\varepsilon}}
\newunicodechar{δ}{\ensuremath{\delta}}
\newunicodechar{π}{\ensuremath{\pi}}
\newunicodechar{σ}{\ensuremath{\sigma}}
\newunicodechar{φ}{\ensuremath{\varphi}}
\newunicodechar{∞}{\ensuremath{\infty}}
\newunicodechar{⌊}{\ensuremath{\lfloor}}
\newunicodechar{⌋}{\ensuremath{\rfloor}}
\newunicodechar{⌈}{\ensuremath{\lceil}}
\newunicodechar{⌉}{\ensuremath{\rceil}}
\newunicodechar{∣}{\ensuremath{\mid}}
\newunicodechar{≡}{\ensuremath{\equiv}}
\newunicodechar{×}{\ensuremath{\times}}
\newunicodechar{⧸}{\ensuremath{/}}
\newunicodechar{‖}{\ensuremath{\|}}
\newunicodechar{√}{\ensuremath{\surd}}
\newunicodechar{∅}{\ensuremath{\emptyset}}
\newunicodechar{⦃}{\ensuremath{\{\!|}}
\newunicodechar{⦄}{\ensuremath{|\!\}}}
\newunicodechar{𝓝}{\ensuremath{\mathcal{N}}}
\newunicodechar{≫}{\ensuremath{\gg}}
\newunicodechar{≪}{\ensuremath{\ll}}

\theoremstyle{plain}
\newtheorem{theorem}{Theorem}
\newtheorem{lemma}[theorem]{Lemma}
\newtheorem{proposition}[theorem]{Proposition}
\newtheorem{corollary}[theorem]{Corollary}
\newtheorem{conjecture}[theorem]{Conjecture}
\theoremstyle{definition}
\newtheorem{definition}[theorem]{Definition}
\newtheorem{remark}[theorem]{Remark}

\emergencystretch=3em
% Lean identifiers in running text (breakable at dots and underscores).
\expandafter\def\expandafter\UrlBreaks\expandafter{\UrlBreaks\do\_\do\.}
\DeclareUrlCommand\decl{\urlstyle{tt}}
\newcommand{\N}{\mathbb{N}}
\newcommand{\Cyc}{\Lambda}
\newcommand{\Lm}[1]{\mathtt{#1}}

\title[The digit-power iteration from 119 ends at 371 (OEIS A115026)]{The digit-power iteration started at $119$ ends at the fixed point $371$: a counterexample to the conjecture in OEIS A115026}
\author{math0}
\address{Produced by Claude (an AI model developed by Anthropic) in the math0 project}
\date{October 18, 2026}

\begin{document}
\begin{abstract}
Let $F(m)$ be the sum of the $L$-th powers of the decimal digits of $m$, where $L$ is the number of digits of $m$ (OEIS A101337), and let $a(n)$ be the smallest element of the cycle that the orbit $n,F(n),F(F(n)),\dots$ eventually enters (OEIS A115026).
In 2022 Harvey P.~Dale conjectured in the OEIS entry that no term $a(n)$, $n\ge1$, is greater than $370$.
The conjecture is false: $119\mapsto731\mapsto371\mapsto371$, so $a(119)=371$.
This counterexample is not new, since a 2009 table of cycles of the same iteration by Hans Havermann in OEIS A151543 already records that $119$ is the smallest number whose orbit ends at the fixed point $371$; we therefore classify this note as the first formal proof of a known result.
The refutation is verified in Lean~4 against a statement that math0 formalized from the OEIS entry, which was reviewed for fidelity and frozen before the proof attempt, and in which $a$ is a computable function proved to return the minimum of the limit cycle and checked by the kernel against all $65$ listed terms.
\end{abstract}
\maketitle

\section{Introduction}

For an integer $m\ge1$ with $L$ decimal digits $d_{L-1}\cdots d_1d_0$, let
\[
  F(m)=\sum_{i=0}^{L-1} d_i^{\,L}.
\]
This is OEIS A101337 \cite{A101337} (Gordon Hamilton, 2004). For example, $F(89)=8^2+9^2=145$ and $F(145)=1^3+4^3+5^3=190$; note that the exponent is the number of digits of the current argument, so it can change along an orbit.
The fixed points of $F$ are the positive narcissistic (Armstrong, pluperfect) numbers $1,\dots,9,153,370,371,407,1634,\dots$ of A005188 \cite{A005188}.
Every orbit of $F$ is eventually periodic, because $F(m)<m$ as soon as $m$ has sufficiently many digits, so every orbit is bounded (Lemma~\ref{lem:bound} below).

The OEIS entry A115026 \cite{A115026}, submitted by Sergio Pimentel on February~24, 2006, is the ``Limiting value of $n$ under iteration'' of $F$: the orbit of $n$ ends in a fixed point or in a cycle of length greater than one, and $a(n)$ is the fixed point, respectively the smallest element of the cycle.
For instance $89\mapsto145\mapsto190\mapsto730\mapsto370\mapsto370$, so $a(89)=370$, and the orbit of $59$ ends in the $3$-cycle $(160,217,352)$, so $a(59)=160$.
On June~8, 2022, Harvey P.~Dale added the comment ``Conjecture: no term is greater than 370'', together with a Mathematica program that computes $a(1),\dots,a(70)$, following each orbit for $20$ steps.
The $65$ listed terms are all at most $370$.
We record that the conjecture is false (Theorem~\ref{thm:main}):
\[
  F(119)=1^3+1^3+9^3=731,\qquad F(731)=7^3+3^3+1^3=371,\qquad F(371)=3^3+7^3+1^3=371,
\]
so the orbit of $119$ ends at the fixed point $371$, and $a(119)=371>370$.

\subsection*{Prior work}
The counterexample was in the OEIS thirteen years before the conjecture was stated.
The same iteration was studied in 2009 under the name ``Recurring Digital Invariant Variant'' (A151543 \cite{A151543} attributes the name to Mensanator on the rec.puzzles site).
The entry A151543, ``Smallest precursor of $n$-th cycle in the `Recurring Digital Invariant Variant' problem'', was created by N.~J.~A. Sloane on May~15, 2009, based on an email from Eric Angelini of February~18, 2009 (see also Angelini's web page \cite{Angelini}), and contains a table, computed by Hans Havermann, of the cycles of $F$ ordered by the smallest number whose orbit ends in them (Havermann's reply to the SeqFan mailing list of February~18, 2009, reproduced on Angelini's page, already gives the first $34$ rows of this table, including the row quoted next).
Its entry ``13 \{ 119, 1, \{371\}\}'' states that $119$ is the smallest number whose orbit ends in the fixed point $371$; in the notation of A115026 this is $a(119)=371$.
The table also lists the fixed points $407$, $1634$, $8208$, \dots{} and cycles of length greater than one with minima $2929$, $18829$, $58618$, \dots, all of which are values of $a$ greater than $370$.
The table moreover implies that $119$ is the smallest counterexample: the cycles reached from some $n<119$ are the first twelve of the table, whose minima are $1,\dots,9$, $370$, $160$ and $153$.
(We confirmed by a direct computation, which is not part of the formal proof, that $a(n)\le370$ for $1\le n\le118$.)
The cycle elements are listed in A151544 \cite{A151544} (Sloane, 2009), and the elements of cycles of length greater than one in A157714 \cite{A157714} (Havermann, March~4, 2009).
Even more directly, $371$ and $407$ are narcissistic numbers, listed in A005188 \cite{A005188} long before 2022; for a fixed point $m$ of $F$ one has $a(m)=m$, so $a(371)=371$ and $a(407)=407$ also contradict the conjecture.

The entry A115026 (revision \#12 of June~8, 2022, the version we read) does not cross-reference A151543 or A005188, has no correction of Dale's comment, and still states the conjecture.
The math0 referees found no published retraction of the comment, and the statement is not among the OEIS statements of the formal-conjectures repository \cite{FC}.
The mathematics is elementary.
What we add is a formal statement of the conjecture, with a definition of $a$ that is proved to agree with the OEIS definition, and a kernel-checked refutation.
We classify the result as the first formal proof of a known result.

\section{Statement}\label{sec:statement}

\subsection*{The OEIS definition and conjecture}
The name of A115026 and its first comment read, verbatim:
\begin{quote}\raggedright\frenchspacing
{\small\texttt{Limiting value of n under iteration of "sum of the digits raised to the power of the number of digits of n" (A101337).}}\\[3pt]
{\small\texttt{Iterate A101337 starting at n until reaching a constant value (like 370) or a cycle (like 160, 217, 352, 160, ...). In the latter case, a(n) takes the smallest value in the cycle (e.g., a(59) = 160). Since k*9\textasciicircum{}k < 10\textasciicircum{}k for all k > 34, each number n is guaranteed to yield a smaller number a(n) if n > 10\textasciicircum{}34, so every number reaches a constant or a cycle under this sequence.}}
\end{quote}
(Its conclusion is correct, but the inequality quoted does not by itself give $F(n)<n$; for $n$ with $L$ digits the bound $F(n)\le L\cdot9^L$ gives $F(n)<n$ once $L\cdot9^L<10^{L-1}$, which holds for all $L\ge61$. We use the threshold $70$ in Lemma~\ref{lem:bound} below.)
The offset is $1$. The second comment is the conjecture, verbatim:
\begin{quote}\raggedright\frenchspacing
{\small\texttt{Conjecture: no term is greater than 370. - \_Harvey P. Dale\_, Jun~08~2022}}
\end{quote}

\begin{definition}\label{def:main}
Let $F\colon\N\to\N$ be as in the introduction, with $F(0)=0$. For $n\in\N$ let
\[
  \Cyc(n)=\{\,m\in\N:\ m=F^k(n)\text{ for some }k\ge0,\ \text{and}\ F^p(m)=m\text{ for some }p\ge1\,\},
\]
the set of periodic points of $F$ on the orbit of $n$, and let $a(n)=\min\Cyc(n)$.
\end{definition}

Call $m$ periodic with period $p\ge1$ if $F^p(m)=m$. If $F^i(n)$ has period $p$, then every $F^j(n)$ with $j\ge i$ is periodic and lies in the cycle $\{F^i(n),\dots,F^{i+p-1}(n)\}$. Conversely, if $m=F^j(n)$ has period $q$, choose $s\ge1$ with $qs\ge i$; then
\[
  m=F^{qs}(m)=F^{qs+j-i}\bigl(F^i(n)\bigr)
\]
lies in that cycle. So $\Cyc(n)$ is exactly the cycle in which the orbit of $n$ ends; a fixed point (``constant value'') is a cycle of length one. The set $\Cyc(n)$ is nonempty by Lemma~\ref{lem:bound} and Proposition~\ref{prop:spec} below.
For $m\ge1$ the exponent $L$ is at least $1$, so a digit $0$ contributes $0^L=0$, and $F(m)\ge1$ because the leading digit is nonzero; hence the value $F(0)=0$ never occurs on the orbit of an $n\ge1$.

\begin{theorem}\label{thm:main}
$a(119)=371$. In particular, the conjecture $a(n)\le370$ for all $n\ge1$ is false.
\end{theorem}

\subsection*{Lean formalization}
The statement was formalized by math0 in the file \path{lean/Math0/Statements/OEIS/A115026.lean}; it was reviewed for fidelity and frozen (SHA-256 \texttt{39ccb0d4\ldots 799b1}) before the proof attempt. The definitions read as follows (verbatim, except that the copyright header, the import, the module docstring, all docstrings, and the theorems stated between \decl{a} and \decl{statement} are omitted):
\begin{lstlisting}
namespace Math0.Statements.A115026
def F (m : ℕ) : ℕ := ((Nat.digits 10 m).map (· ^ (Nat.digits 10 m).length)).sum
def limitCycle (n : ℕ) : Set ℕ :=
  {m | (∃ k, F^[k] n = m) ∧ ∃ p, 0 < p ∧ F^[p] m = m}
def meet : ℕ → ℕ → ℕ → ℕ
  | 0, x, _ => x
  | fuel + 1, x, z => if x = z then x else meet fuel (F x) (F (F z))
def cycleMin (y : ℕ) : ℕ → ℕ → ℕ → ℕ
  | 0, _, m => m
  | fuel + 1, cur, m => if cur = y then m else cycleMin y fuel (F cur) (min m cur)
def C : ℕ := 70 * 9 ^ 70
def fuel (n : ℕ) : ℕ := (max n C + 2) ^ 2
def a (n : ℕ) : ℕ :=
  let y := meet (fuel n) (F n) (F (F n))
  cycleMin y (fuel n) (F y) y
def statement : Prop :=
  ∀ n : ℕ, 1 ≤ n → a n ≤ 370
end Math0.Statements.A115026
\end{lstlisting}
Here \texttt{Nat.digits 10 m} is the list of decimal digits of $m$, least significant first (empty for $m=0$), so \decl{F} is $F$, and \decl{limitCycle} is $\Cyc$.
The function \decl{a} is computable, so that the kernel can evaluate it: \decl{meet} is Floyd's tortoise-and-hare loop, which compares $F^k(n)$ with $F^{2k}(n)$ for $k=1,2,\dots$ and returns the first $F^k(n)$ with $F^k(n)=F^{2k}(n)$, and \decl{cycleMin} walks around the cycle of the returned point $y$ and keeps the minimum. Both loops carry an explicit fuel argument.
Among other lemmas, the file proves the following (statements verbatim, proofs omitted):
\begin{lstlisting}
theorem a_data : (List.range 65).map (fun i => a (i + 1)) =
    [1, 2, 3, 4, 5, 6, 7, 8, 9, 1, 2, 5, 1, 370, 370, 370, 370, 370, 1, 4, 5, 8, 1, 4, 370, 370,
      370, 1, 370, 9, 1, 1, 370, 370, 370, 370, 370, 370, 370, 370, 370, 4, 370, 1, 370, 370,
      370, 370, 1, 370, 370, 370, 370, 370, 370, 370, 370, 370, 160, 370, 370, 370, 370, 370,
      370]
theorem a_spec (n : ℕ) : a n ∈ limitCycle n ∧ ∀ m ∈ limitCycle n, a n ≤ m
\end{lstlisting}
The theorem \decl{a_data}, proved by \texttt{decide +kernel}, checks that $a(1),\dots,a(65)$ are the $65$ terms listed in the OEIS (the full \texttt{\%S}, \texttt{\%T}, \texttt{\%U} lines of the entry).
The theorem \decl{a_spec} states that \decl{a} computes $\min\Cyc(n)$ for every $n$; so the Lean function \decl{a} is the sequence of Definition~\ref{def:main}, not a conjectured formula for it.
Since the offset is $1$, ``no term is greater than $370$'' is $\forall n\ge1,\ a(n)\le370$, which is \decl{statement}; it adds no hypotheses.
We sketch the proof of \decl{a_spec}, following the lemmas of the statement file.

\begin{lemma}[\decl{F_le}, \decl{key_ineq}, \decl{F_le_max}, \decl{iterate_le}]\label{lem:bound}
Let $C=70\cdot9^{70}$. For all $m,n\in\N$ and $k\ge0$ we have $F(m)\le\max(m,C)$ and $F^k(n)\le\max(n,C)$.
\end{lemma}

\begin{proof}
If $m$ has $L$ digits, each at most $9$, then $F(m)\le L\cdot9^L$. If $L<70$ this is at most $C$. If $L\ge70$, then $10L\cdot9^L<10^L$ (by induction on $L$: it holds for $L=70$, and $10(L+1)9^{L+1}=9\cdot10(L+1)9^L\le10\cdot10L\cdot9^L<10^{L+1}$ since $9(L+1)\le10L$), and $10^L\le10m$ because $m\ge10^{L-1}$; hence $10F(m)<10m$. The bound for $F^k(n)$ follows by induction on $k$.
\end{proof}

\begin{proposition}[\decl{a_spec}]\label{prop:spec}
For every $n\in\N$, $a(n)\in\Cyc(n)$ and $a(n)\le m$ for every $m\in\Cyc(n)$.
\end{proposition}

\begin{proof}[Sketch]
Let $M=\max(n,C)$, so that \texttt{fuel n} is $(M+2)^2$. By Lemma~\ref{lem:bound} the $M+2$ values $F^0(n),\dots,F^{M+1}(n)$ lie in $\{0,\dots,M\}$, so by the pigeonhole principle $F^i(n)=F^j(n)$ for some $i<j\le M+1$ (\decl{exists_meet}). Then $F^i(n)$ has period $j-i$, and $K=(j-i)(i+1)$ satisfies $1\le K\le(M+2)^2$, $K\ge i$ and $F^K(n)=F^{2K}(n)$ (\decl{meet_of_repeat}). Hence Floyd's loop, which at its $k$-th step compares $F^k(n)$ with $F^{2k}(n)$, stops at some $k\le K$ before its fuel runs out and returns $y=F^k(n)$ with $F^k(y)=y$ (\decl{meet_spec}). The walk \decl{cycleMin} visits $F^1(y),F^2(y),\dots$ until it returns to $y$, which happens after $t\le k$ steps, and returns $\min\{F^s(y):0\le s<t\}$, which is the minimum of the whole orbit of $y$ because $F^s(y)=F^{s\bmod t}(y)$ (\decl{cycleMin_spec}). This value is $F^{s+k}(n)$ for some $s$ and is periodic, so it lies in $\Cyc(n)$. Conversely, if $m=F^i(n)$ satisfies $F^q(m)=m$ with $q\ge1$, then $m=F^{qk}(m)=F^{qk+i-k}(y)$ lies on the orbit of $y$, so $a(n)\le m$.
\end{proof}

\section{Proof}\label{sec:proof}

The Lean proof evaluates the function \decl{a} of Section~\ref{sec:statement} at $119$ in the kernel. We give the corresponding computation by hand, and then check the result against Definition~\ref{def:main} directly.

\begin{lemma}\label{lem:orbit}
$F(119)=731$, $F(731)=371$ and $F(371)=371$. Hence $F^k(119)=371$ for all $k\ge2$.
\end{lemma}

\begin{proof}
All three numbers have three digits, so the exponent is $3$: $1+1+729=731$, $343+27+1=371$ and $27+343+1=371$.
\end{proof}

\begin{lemma}[\decl{a_119}]\label{lem:a119}
The function \decl{a} of Section~\ref{sec:statement} satisfies $a(119)=371$.
\end{lemma}

\begin{proof}
Let $N=\Lm{fuel}(119)=(C+2)^2$, where $C=70\cdot9^{70}>119$; all that matters below is $N\ge2$. By definition, $a(119)=\Lm{cycleMin}(y,N,F(y),y)$ with $y=\Lm{meet}(N,F(119),F(F(119)))$. By Lemma~\ref{lem:orbit},
\[
  y=\Lm{meet}(N,731,371)=\Lm{meet}(N-1,F(731),F(F(371)))=\Lm{meet}(N-1,371,371)=371,
\]
where the first step uses $731\ne371$ and $N\ge1$, and the last uses $371=371$ and $N-1\ge1$. Then
\[
  a(119)=\Lm{cycleMin}(371,N,F(371),371)=\Lm{cycleMin}(371,N,371,371)=371,
\]
because $N\ge1$ and the current point $371$ equals $y$, so the loop returns its accumulator $371$ at once.
\end{proof}

\begin{proof}[Proof of Theorem~\ref{thm:main}]
By Lemma~\ref{lem:a119} and Proposition~\ref{prop:spec}, $a(119)=371$ for the sequence of Definition~\ref{def:main}. If $a(n)\le370$ held for all $n\ge1$, then in particular $371=a(119)\le370$, which is absurd. This is the Lean theorem \decl{result}: it applies the hypothesis at $n=119$, rewrites with $a(119)=371$, and derives a contradiction from the hypothesis $371\le370$ with \texttt{omega}.
\end{proof}

\begin{remark}[direct check]\label{rem:direct}
The value $a(119)=371$ can also be read off from Definition~\ref{def:main} without the algorithm. By Lemma~\ref{lem:orbit} the orbit of $119$ is $\{119,731,371\}$. The point $371$ is periodic with period $1$. The point $119$ is not periodic, because $F^p(119)\in\{731,371\}$ for all $p\ge1$, and $731$ is not periodic, because $F^p(731)=371$ for all $p\ge1$. So $\Cyc(119)=\{371\}$ and $a(119)=371$. In the Lean development this argument is not needed, since Proposition~\ref{prop:spec} holds for every $n$.
\end{remark}

\begin{remark}[the certificate]\label{rem:kernel}
The Lean proof of Lemma~\ref{lem:a119} is \texttt{decide +kernel}: the kernel reduces the decision procedure for the equality \texttt{a 119 = 371} on natural numbers to \texttt{true}, unfolding \decl{a}, \decl{fuel}, \decl{meet}, \decl{cycleMin}, \decl{F} and \texttt{Nat.digits}. Arithmetic on numerals, including $C\approx4.4\cdot10^{68}$ and $N\approx1.9\cdot10^{137}$, uses the kernel's built-in arbitrary-precision natural numbers, and only the two iterations of \decl{meet} and the single iteration of \decl{cycleMin} described in the proof of Lemma~\ref{lem:a119} are executed, so the size of the fuel is harmless. The theorem \decl{a_data} of the statement file is checked in the same way. Unlike \texttt{native\_decide}, this introduces no axiom and does not trust the compiler.
\end{remark}

\begin{remark}[further counterexamples; not formalized]
By direct computation, the integers $n\le200$ with $a(n)>370$ begin $119,137,147,159,168,173,\dots$, with $a(137)=a(173)=371$, $a(147)=18829$ (the minimum of a cycle of length $14$) and $a(177)=58618$ (the minimum of a cycle of length $2$). Every narcissistic number $m\ge371$, such as $371$, $407$ or $1634$, satisfies $a(m)=m$. The Lean proof uses only $n=119$.
\end{remark}

\section{Machine verification}\label{sec:lean}

\subsection*{Files and declarations}
The proof is in the file \path{lean/Math0/Frontier/OEIS/A115026.lean} of the math0 repository, \url{https://github.com/Foundation-0/math0}. It imports the frozen statement file \path{lean/Math0/Statements/OEIS/A115026.lean} (Section~\ref{sec:statement}) and the audit commands in \path{lean/Math0/Audit.lean}. The main declaration is \decl{Math0.Frontier.A115026.result}, of type $\neg$\,\decl{Math0.Statements.A115026.statement}. The file reads, verbatim (the copyright header, imports, module docstring and docstrings omitted):
\begin{lstlisting}
namespace Math0.Frontier.A115026
open Math0.Statements.A115026
private theorem a_119 : a 119 = 371 := by
  decide +kernel
theorem result : ¬ statement := by
  intro h
  have h119 := h 119 (by norm_num)
  rw [a_119] at h119
  omega
end Math0.Frontier.A115026
#math0_settles Math0.Frontier.A115026.result Math0.Statements.A115026.statement
\end{lstlisting}
The file uses no \texttt{native\_decide}, \texttt{sorry}, \texttt{axiom} or \texttt{set\_option}. It does not use \decl{a_spec}: the theorem \decl{result} refutes \decl{statement} as stated, about the function \decl{a}, and \decl{a_spec} is what identifies that function with the OEIS sequence.

\subsection*{Audit}
The command \texttt{\#math0\_settles}, defined in \path{lean/Math0/Audit.lean} for statements that math0 formalizes itself, fails unless (1) \decl{statement} is a \texttt{Prop}-valued definition without universe parameters; (2) \decl{statement} and \decl{result} use no axioms besides \decl{propext}, \decl{Classical.choice} and \decl{Quot.sound} (this rejects \decl{sorryAx}, the axiom \decl{Lean.ofReduceBool} introduced by \texttt{native\_decide}, and custom axioms); and (3) the type of \decl{result} is syntactically the statement or its negation, up to metadata annotations, with no unfolding. On this file it reports
\begin{lstlisting}[basicstyle=\ttfamily\scriptsize]
math0 audit: 'Math0.Statements.A115026.statement' depends on axioms [Classical.choice, Quot.sound, propext]
math0 audit: 'Math0.Frontier.A115026.result' depends on axioms [Classical.choice, Quot.sound, propext]
math0 statement check: 'Math0.Frontier.A115026.result' refutes 'Math0.Statements.A115026.statement'
\end{lstlisting}
so \decl{result} proves the negation of the statement and depends only on the three standard axioms. The theorems \decl{a_data} and \decl{a_spec} of the statement file also depend only on these axioms.

\subsection*{Toolchain and reproduction}
We use Lean~4 \cite{Lean4} \texttt{v4.33.1} and Mathlib \cite{Mathlib} \texttt{v4.33.1}, pinned in \path{lean/}; the repository also pins formal-conjectures \cite{FC} at commit \texttt{c77fdd4}, which this result does not use.
To reproduce, run \path{scripts/verify.sh} from the repository root: it fetches the Mathlib build cache, builds all math0 results with \texttt{lake build} (re-running every audit command), checks that the pinned \texttt{formal\_conjectures} and \texttt{mathlib} checkouts are unmodified, and reports the number of passed statement checks. The single file can be checked with \texttt{lake env lean Math0/Frontier/OEIS/A115026.lean} in the directory \path{lean}; this takes about a minute or less on our machine, most of it spent loading Mathlib.

{\small
\begin{thebibliography}{99}

\bibitem{Angelini}
E.~Angelini, \emph{A Recurring Digital Invariant Variant}, web page, \url{http://www.cetteadressecomportecinquantesignes.com/RecurDigit.htm}; cached copy at \url{https://oeis.org/A151543/a151543.pdf}.

\bibitem{FC}
M.~Firsching, P.~Lezeau, S.~Mercuri, M.~Z. Horv\'ath, Y.~Dillies, C.~S\"onne, E.~Wieser, F.~Zhang, T.~Hubert, B.~Ag\"uera y Arcas, and P.~Kohli, \emph{Formal Conjectures: An Open and Evolving Benchmark for Verified Discovery in Mathematics}, arXiv:2605.13171, 2026. Repository: \url{https://github.com/google-deepmind/formal-conjectures}.

\bibitem{Mathlib}
The mathlib Community, \emph{The Lean mathematical library}, in: Proceedings of the 9th ACM SIGPLAN International Conference on Certified Programs and Proofs (CPP 2020), ACM, 2020, 367--381.

\bibitem{Lean4}
L.~de Moura and S.~Ullrich, \emph{The Lean 4 theorem prover and programming language}, in: Automated Deduction -- CADE 28, Lecture Notes in Comput. Sci. \textbf{12699}, Springer, 2021, 625--635.

\bibitem{A005188}
The OEIS Foundation Inc., \url{https://oeis.org/A005188} (Armstrong, or narcissistic, numbers), entry by N.~J.~A. Sloane and R.~G. Wilson~v.

\bibitem{A101337}
The OEIS Foundation Inc., \url{https://oeis.org/A101337}, entry by G.~Hamilton (Dec.~24, 2004).

\bibitem{A115026}
The OEIS Foundation Inc., \url{https://oeis.org/A115026}, entry by S.~Pimentel (Feb.~24, 2006); conjecture by H.~P. Dale (June~8, 2022); revision \#12, June~8, 2022.

\bibitem{A151543}
The OEIS Foundation Inc., \url{https://oeis.org/A151543}, entry by N.~J.~A. Sloane (May~15, 2009), based on an email from E.~Angelini (Feb.~18, 2009), with a table of cycles by H.~Havermann.

\bibitem{A151544}
The OEIS Foundation Inc., \url{https://oeis.org/A151544}, entry by N.~J.~A. Sloane (May~15, 2009), based on an email from E.~Angelini.

\bibitem{A157714}
The OEIS Foundation Inc., \url{https://oeis.org/A157714} (base-10 pseudo-altruistic numbers), entry by H.~Havermann (Mar.~4, 2009).

\end{thebibliography}
}

\end{document}
